\documentclass[a4paper, serif, 10pt]{moderncv}

%% moderncv
% style and color
\moderncvstyle{banking}
\moderncvcolor{black}

%% page layout/margins
\usepackage[top=2.5cm, bottom=2.5cm, left=1.5cm, right=1.5cm]{geometry}

\nopagenumbers{}

%% Babel config for Norwegian language
% ref:
% - https://latex3.github.io/babel/guides/locale-norwegian.html
\usepackage[norsk]{babel}

%% fontspec
\usepackage{fontspec}

\IfFontExistsTF{Linux Libertine}{
  \setmainfont[Ligatures={TeX, Common}, Numbers=OldStyle]{Linux Libertine}
}{}
\IfFontExistsTF{Linux Libertine O}{
  \setmainfont[Ligatures={TeX, Common}, Numbers=OldStyle]{Linux Libertine O}
}{}

%% CTeX
% CJK support, used to show CJK characters in the resume
%
% - fontset=none: disable builtin fontset but instead set the CJK font manually
% - heading=false: disable ctex heading
% - punct=kaiming: use kaiming punctuations styles for CJK
% - scheme=plain: use plain scheme, do not override \normalsize font size
% - space=auto: space settings for CJK characters
%
% ref:
% - http://ctan.mirrorcatalogs.com/language/chinese/ctex/ctex.pdf
\usepackage[UTF8, heading=false, punct=kaiming, scheme=plain, space=auto]{ctex}

\IfFontExistsTF{Noto Serif CJK SC}{
  \setCJKmainfont{Noto Serif CJK SC}
}{}
\IfFontExistsTF{Noto Sans CJK SC}{
  \setCJKsansfont{Noto Sans CJK SC}
}{}

%% line spacing
\usepackage{setspace}
\setstretch{1.125}

%% URL styling - use normal text instead of monospace
\urlstyle{same}

% Auto-underline all links
% Defer until \begin{document} so hyperref is already loaded
\AtBeginDocument{%
  \let\oldhref\href
  \renewcommand{\href}[2]{\oldhref{#1}{\underline{#2}}}%
}

%% Patch \httplink and \httpslink to handle full URLs
% The original moderncv macros always prepend http:// or https:// to URLs.
% This causes issues when the URL already has a protocol (e.g., https://example.com)
% resulting in malformed URLs like https://https://example.com.
% This patch checks if the URL already contains "://" and uses it directly if so.
% Note: We use \str_set:Nx (x-expansion) to expand macros like \@homepage first.
\ExplSyntaxOn
\str_new:N \g_yamlresume_url_str

\RenewDocumentCommand{\httpslink}{O{}m}{
  \str_set:Nx \g_yamlresume_url_str {#2}
  \str_if_in:NnTF \g_yamlresume_url_str {://}
    { \url{#2} }
    { \url{https://#2} }
}

\RenewDocumentCommand{\httplink}{O{}m}{
  \str_set:Nx \g_yamlresume_url_str {#2}
  \str_if_in:NnTF \g_yamlresume_url_str {://}
    { \url{#2} }
    { \url{http://#2} }
}
\ExplSyntaxOff

\name{Ingrid Solberg}{}
\title{Erfaren prosjektleder med lidenskap for smidig metodikk og verdiskapning}
\phone[mobile]{+47 123 45 678}
\email{ingrid.solberg@example.no}
\homepage{https://www.ingridsolberg.no}

\address{Karl Johans gate 15, Oslo, Oslo, Norge, 0150}{}{}

\extrainfo{{\small \faLinkedin}\ \href{https://www.linkedin.com/in/ingridsolberg}{@ingridsolberg} {} {} {} • {} {} {} 
{\small \faGithub}\ \href{https://github.com/ingridsolberg}{@ingridsolberg} {} {} {} • {} {} {} 
{\small \faTwitter}\ \href{https://twitter.com/ingridsolberg}{@ingridsolberg}}

\begin{document}

\maketitle

\section{Grunnleggende informasjon}

\cvline{}{\begin{itemize}
\item Prosjektleder med over 8 års erfaring fra teknologi- og energibransjen
\item Dokumentert suksess med å lede tverrfaglige team og levere prosjekter innen tid og budsjett
\item Sterk kompetanse innen smidige metoder (Scrum, Kanban) og tradisjonell prosjektstyring (PMBOK)
\item Strategisk tenker med fokus på kontinuerlig forbedring og bærekraftige løsninger
\item Gode kommunikasjons- og lederegenskaper, vant til å samarbeide med interessenter på alle nivåer
\end{itemize}}
\section{Utdanning}

\cventry{aug. 2015–juni 2017}
        {Master, Industriell økonomi og teknologiledelse, Poeng: 3.8}
        {Norges teknisk-naturvitenskapelige universitet (NTNU)}
        {\url{https://www.ntnu.no/}}
        {}
        {\begin{itemize}
\item Spesialisering i prosjektledelse og innovasjonsledelse
\item Gjennomførte masteroppgave om smidig prosjektstyring i energisektoren
\item Verv som leder for studentorganisasjonen Prosjektforum
\end{itemize}
\textbf{Kurs}: Prosjektledelse, Risikostyring, Smidig utvikling, Strategisk ledelse, Finansstyring, Forhandlingsteknikk}

\cventry{aug. 2012–juni 2015}
        {Bachelor, Økonomi og administrasjon, Poeng: 3.5}
        {Handelshøyskolen BI}
        {\url{https://www.bi.no/}}
        {}
        {\begin{itemize}
\item Grunnleggende forståelse for forretningsdrift og økonomistyring
\item Deltok i case-konkurranser og studentprosjekter
\end{itemize}
\textbf{Kurs}: Bedriftsøkonomi, Markedsføring, Organisasjonspsykologi, Statistikk}
\section{Arbeidserfaring}

\cventry{jan. 2021–Nå}
        {Senior Prosjektleder}
        {Equinor}
        {\url{https://www.equinor.com/}}
        {}
        {\begin{itemize}
\item Leder komplekse tverrfaglige prosjekter innen fornybar energi og digitalisering
\item Ansvarlig for prosjektbudsjetter på opptil 200 MNOK og team på 20+ medlemmer
\item Implementerte smidige arbeidsmetoder som reduserte leveringstiden med 25\%
\item Bygget sterke relasjoner med interessenter og leverandører for å sikre fremdrift
\item Veileder og mentor for junior prosjektledere i organisasjonen
\end{itemize}
\textbf{Nøkkelord}: Prosjektledelse, Smidig, Fornybar energi, Digitalisering, Risikostyring}

\cventry{sep. 2017–des. 2020}
        {Prosjektleder}
        {Telenor}
        {\url{https://www.telenor.no/}}
        {}
        {\begin{itemize}
\item Ledet flere IT-utviklingsprosjekter fra konsept til produksjon
\item Koordinerte interne team og eksterne leverandører for å levere på tid og budsjett
\item Innførte Kanban for bedre flyt og oversikt i utviklingsteamene
\item Gjennomførte risikovurderinger og utarbeidet beredskapsplaner
\end{itemize}
\textbf{Nøkkelord}: IT-prosjekter, Kanban, Leverandørstyring, Agile}

\cventry{juli 2016–aug. 2017}
        {Prosjektkoordinator}
        {DNV}
        {\url{https://www.dnv.no/}}
        {}
        {\begin{itemize}
\item Støttet prosjektledere i planlegging, oppfølging og rapportering
\item Utviklet maler og verktøy for prosjektstyring som ble brukt på tvers av avdelingen
\item Bidro til å etablere et prosjektkontor (PMO) for bedre styring og kontroll
\end{itemize}
\textbf{Nøkkelord}: Prosjektstøtte, PMO, Rapportering}
\section{Språk}

\cvline{Norsk}{Morsmål eller tospråklig nivå \hfill \textbf{Nøkkelord}: Morsmål}
\cvline{Engelsk}{Fullt profesjonelt nivå \hfill \textbf{Nøkkelord}: TOEFL 110, Forretningsflyt}
\section{Ferdigheter}

\cvline{Prosjektledelse}{Ekspert \hfill \textbf{Nøkkelord}: PMBOK, PRINCE2, Risikostyring, Budsjettering, Interessentanalyse}
\cvline{Smidig metodikk}{Avansert \hfill \textbf{Nøkkelord}: Scrum, Kanban, Lean, SAFe}
\cvline{Verktøy}{Avansert \hfill \textbf{Nøkkelord}: Jira, Microsoft Project, Trello, Asana, Miro}
\cvline{Ledelse}{Ekspert \hfill \textbf{Nøkkelord}: Teamledelse, Mentoring, Kommunikasjon, Forhandling}
\section{Utmerkelser}

\cventry{nov. 2022}
        {Norsk Prosjektlederforening}
        {Årets Prosjektleder}
        {}
        {}
        {\begin{itemize}
\item Tildelt for fremragende ledelse av digitaliseringsprosjektet "Fremtidens Energi"
\item Anerkjent for innovativ bruk av smidige metoder og sterk resultatorientering
\end{itemize}}

\cventry{juni 2017}
        {NTNU}
        {Beste Masteroppgave}
        {}
        {}
        {\begin{itemize}
\item Tildelt for masteroppgaven "Smidig prosjektstyring i energisektoren"
\item Karakter A og ros for praktisk relevans
\end{itemize}}
\section{Sertifikater}

\cventry{mars 2019}
        {Project Management Institute}
        {Project Management Professional (PMP)}
        {\url{https://www.pmi.org/certifications/project-management-pmp}}
        {}
        {}

\cventry{juni 2016}
        {AXELOS}
        {PRINCE2 Foundation}
        {\url{https://www.axelos.com/certifications/prince2}}
        {}
        {}
\section{Publikasjoner}

\cventry{sep. 2017}
        {Smidig prosjektstyring i energisektoren}
        {Magma - Econa}
        {\url{https://www.magma.no/smidig-prosjektstyring}}
        {}
        {\begin{itemize}
\item Artikkel basert på masteroppgave om hvordan smidige metoder kan forbedre prosjektgjennomføring i energibransjen
\item Presenterer case-studier og anbefalinger for ledere
\end{itemize}}
\section{Prosjekter}

\cventry{feb. 2021–des. 2022}
        {Digital plattform for sanntidsstyring av fornybare energikilder}
        {Fremtidens Energi}
        {\url{https://www.equinor.com/fremtidens-energi}}
        {}
        {\begin{itemize}
\item Ledet et team på 15 utviklere, designere og data scientists
\item Leverte plattformen 2 måneder før planlagt og 10\% under budsjett
\item Implementerte smidige metoder og automatisert rapportering
\end{itemize}
\textbf{Nøkkelord}: Fornybar energi, Digitalisering, Smidig, Prosjektledelse}
\section{Interesser}

\cvline{Friluftsliv}{Fjellturer, Ski, Kajakk}
\cvline{Teknologi}{Kunstig intelligens, Bærekraft, Åpen kildekode}
\section{Frivillig arbeid}

\cventry{jan. 2020–Nå}
        {Mentor}
        {Norsk Prosjektlederforening}
        {\url{https://www.prosjektlederforeningen.no/}}
        {}
        {\begin{itemize}
\item Mentorer unge prosjektledere og studenter innen prosjektfaget
\item Holder foredrag og workshops om smidig prosjektledelse
\end{itemize}}
    

\cventry{mars 2018–des. 2019}
        {Rådgiver}
        {StartUp Norge}
        {\url{https://www.startupnorge.no/}}
        {}
        {\begin{itemize}
\item Rådgav oppstartsbedrifter om prosjektplanlegging og finansiering
\item Bidro til å utvikle strukturerte prosjektplaner for tidligfase-selskaper
\end{itemize}}
    
\end{document}